\begin{tikzpicture}[>=latex, font=\rmfamily, scale=1.0]

% ==========================================
% LEFT PANEL: Strut cross-section
% ==========================================
% Center anchor point
\fill[black!60] (0,0) circle (1.5pt);

% Initial state (t=0) - Gives the circles meaning as a "receding" boundary
\draw[thick, gray, dashed] (0,0) circle (3.5);
\node[font=\scriptsize, gray, fill=white, inner sep=2pt] at (0,3.5) {initial boundary ($t=0$)};

% Current gel layer (outer)
\draw[thick, gelcyan, fill=cyan!10] (0,0) circle (2.9);
\node[gelcyan, font=\scriptsize, fill=white, inner sep=1pt] at (1.8,2.7) {gel layer};

% Current unreacted core (inner)
\draw[thick, coregray, fill=gray!25] (0,0) circle (2.2);
\node[coregray, font=\scriptsize, align=center] at (0,1.4) {unreacted \\ core};

% Annotations for thickness and radius
\draw[<->, thick, gelcyan] (2.2,0) -- (2.9,0);
\node[gelcyan, font=\scriptsize, fill=cyan!10, inner sep=1pt] at (2.55,-0.25) {$\delta_g(t)$};

% Radius arrow from center
\draw[->, thick, archgreen] (0,0) -- (1.55, 1.55) node[midway, above left, font=\scriptsize, inner sep=1pt] {$r_s(t)$};

% Receding direction arrow
\draw[->, very thick, rtzred] (3.5,0) -- (2.9,0);
\node[rtzred, font=\scriptsize, fill=white, inner sep=2pt] at (4.1,0.3) {receding};
\node[font=\scriptsize, gray, fill=white, inner sep=2pt] at (4.1,2.0) {pore solution};


% ==========================================
% RIGHT PANEL: Porosity/Evolution Curve
% ==========================================
% Using a scope to reset coordinates to (0,0) for the graph, making it cleaner
\begin{scope}[xshift=6.5cm]
    
    % X-Axis
    \draw[->, thick] (0,0) -- (4.8,0) node[right, font=\scriptsize] {Time, $t$};
    \node[below left, font=\tiny] at (0,0) {0};
    
    % Left Y-axis (Porosity & Diffusivity)
    \draw[->, thick] (0,0) -- (0,3.8) node[above, font=\scriptsize] {$\varepsilon(t), \Veff(t)$};
    
    % Right Y-axis (Tortuosity)
    \draw[->, thick] (4.0,0) -- (4.0,3.8) node[above, font=\scriptsize] {$\bar{\kappa}(t)$};
    
    % Ticks for right axis
    \draw[black!60, thin] (4.0,1.0) -- (4.15,1.0) node[right, font=\tiny] {1};
    \draw[black!60, thin] (4.0,2.0) -- (4.15,2.0) node[right, font=\tiny] {2};
    \draw[black!60, thin] (4.0,3.0) -- (4.15,3.0) node[right, font=\tiny] {3};

    % Curves
    % Porosity \varepsilon(t)
    \draw[very thick, statepurple] (0.1,0.3) .. controls (1.2,1.3) and (2.2,2.2) .. (3.9,3.2);
    \node[statepurple, font=\scriptsize, fill=white, inner sep=2pt] at (3.0,3.2) {$\varepsilon(t)$};

    % Tortuosity \bar{\kappa}(t)
    \draw[very thick, archgreen, dashed] (0.1,2.6) .. controls (1.5,2.2) and (2.7,1.5) .. (3.9,0.8);
    \node[archgreen, font=\scriptsize, fill=white, inner sep=2pt] at (2.8,1 ) {$\bar{\kappa}(t)$};
    
    % Collapse indicator
    \node[font=\scriptsize, gray, fill=white, inner sep=2pt] at (3.3, 0.4) {collapse};
    \draw[->, gray, thin] (3.4, 0.55) -- (3.8, 0.75);

    % Effective Volume/Diffusivity \Veff(t)
    \draw[thick, rtzred] (1.7,0.3) .. controls (2.3,1.6) and (2.9,2.0) .. (3.9,2.2);
    \node[rtzred, font=\scriptsize, fill=white, inner sep=2pt] at (2.7,2) {$\Veff(t)$};

\end{scope}

\end{tikzpicture}
